\section*{Capítulo \ref{ch:b2:climate-proxies} --- Archivos e indicadores climáticos}

\begin{solution}{exo:b2:climate-proxies:1}
(a) $\delta = (27.5 \times (-30) + 1350 \times 0)/(1350 + 27.5) = -825/1377.5 = -0.60$ por mil.
(b) El océano contendría entonces toda el agua: $-0.60$ por mil, \num{0.60} por mil más ligero que
hoy. (c) $27.5/1377.5 = \qty{2.0}{\percent}$.
\end{solution}

\begin{solution}{exo:b2:climate-proxies:2}
(a) $T = 16.1 - 4.76\,(2.50 - (0 - 0.27)) = 16.1 - 4.76 \times 2.77 = 16.1 - 13.19 =
\qty{2.9}{\celsius}$. (b) $T = 16.1 - 4.76\,(2.50 - 0.73) = 16.1 - 8.43 = \qty{7.7}{\celsius}$: el
mismo caparazón da una temperatura casi cinco grados más alta si el agua era más pesada, y por eso
hay que conocer la composición del agua. (c) El caparazón es $2.71 - 2.50 = 0.21$ por mil más pesado,
así que su agua estaba $4.76 \times 0.21 = \qty{1.0}{\celsius}$ más fría, a \qty{1.9}{\celsius}.
\end{solution}

\begin{solution}{exo:b2:climate-proxies:3}
(a) $0.4\,\mathrm{e}^{0.9} = \qty{0.98}{mmol/mol}$ y $0.4\,\mathrm{e}^{1.8} =
\qty{2.42}{mmol/mol}$. (b) El factor es $\mathrm{e}^{0.9} = 2.46$; por grado, $\mathrm{e}^{0.09}
= 1.094$, un crecimiento del \qty{9.4}{\percent} por grado. (c) $T = \ln(\text{Mg/Ca}/0.4)/0.09$:
$\ln 5/0.09 = \qty{17.9}{\celsius}$ y $\ln 2.5/0.09 = \qty{10.2}{\celsius}$.
\end{solution}

\begin{solution}{exo:b2:climate-proxies:4}
(a) $2668 - 2213 + 1 = 456$ años. (b) A un solo árbol puede faltarle localmente un anillo, o puede mostrar un anillo
ilegible; emparejar las secuencias de anillos anchos y estrechos de muchos árboles revela los
anillos que faltan y fija el año de cada anillo, y la concordancia entre árboles mide la
fiabilidad de la datación. (c) En el límite del bosque un factor, el frío de la estación de crecimiento,
limita a todos los árboles, así que comparten la misma secuencia de años buenos y malos; en una llanura
templada y húmeda ningún factor domina, las secuencias de los árboles difieren y la \omterm{def:b2:climate-proxies:dendro}{sincronización} puede fallar.
\end{solution}

\begin{solution}{exo:b2:climate-proxies:5}
(a) $125 \times 40/(3682.2 - 125) = 5000/3557.2 = 1.41$ por mil; $134 \times 40/3548.2 =
5360/3548.2 = 1.51$ por mil. (b) Con $H - h = 3552.2$: $3900/3552.2 = 1.10$, $4550/3552.2 = 1.28$
y $6500/3552.2 = 1.83$ por mil. (c) $\delta_{\mathrm{ice}} = -1.0 \times 3552.2/130 = -27.3$ por
mil, más pesado que cualquier casquete actual (Groenlandia, de $-30$ a $-35$; el casquete de América del Norte,
$-40$; el hielo antártico, hasta $-50$). El balance de masas con un hielo realista y la medida en el agua
intersticial no concuerdan; la diferencia de aproximadamente medio por mil es una incertidumbre real del método.
\end{solution}

\begin{solution}{exo:b2:climate-proxies:6}
(a) $0.372\,\mathrm{e}^{0.0936T} = 0.4\,\mathrm{e}^{0.09T}$ da $\mathrm{e}^{0.0036T} =
0.4/0.372 = 1.0753$, así que $T = \ln 1.0753/0.0036 = \qty{20.2}{\celsius}$. (b) A
\qty{2}{\celsius}: $0.4\,\mathrm{e}^{0.18} = 0.479$ y $0.372\,\mathrm{e}^{0.1872} = 0.449$, así que la
razón ajustada es $0.449/0.479 = 0.937$, un \qty{6.3}{\percent} menor. (c) $T = \ln(0.449/0.4)/0.09 =
\ln 1.1225/0.09 = \qty{1.3}{\celsius}$: la ley predicha lee el caparazón unos
\qty{0.7}{\celsius} demasiado frío. En agua fría, la elección de la calibración importa al nivel de un
grado.
\end{solution}

\begin{solution}{exo:b2:climate-proxies:7}
(a) Del Holoceno a GI-1, \qty{112.19}{m} de hielo para 2989 años: $11\,219/2989 =
\qty{3.75}{cm}$ por año; de GI-16 a GI-17, \qty{17.89}{m} para 1160 años: $1789/1160 =
\qty{1.54}{cm}$ por año. (b) El hielo
fluye: a medida que cada capa queda enterrada se extiende lateralmente y se adelgaza, y las capas más profundas y antiguas
llevan más tiempo adelgazándose. (c) Cuando las capas se vuelven demasiado delgadas, un centímetro o menos, los instrumentos
ya no pueden separar un año del siguiente, y el \omterm{def:b2:climate-proxies:age-model}{modelo de edad} tiene que usar en su lugar puntos de anclaje o modelos de
flujo.
\end{solution}

\begin{solution}{exo:b2:climate-proxies:8}
(a) La fracción del intervalo es $(1550 - 1492.45)/(1604.64 - 1492.45) = 57.55/112.19 = 0.513$,
así que la edad es $11\,703 + 0.513 \times 2989 = 11\,703 + 1533 = \num{13236}$ años b2k. (b) 1950 está 50
años antes de 2000: \num{13186} años antes de 1950. (c) $50 + 0.513 \times (93 - 50) = 72$ años.
\end{solution}

\begin{solution}{exo:b2:climate-proxies:9}
(a) Cada capa dudosa aporta medio año: $2 \times 93 = 186$ y $2 \times 1287 = 2574$
capas dudosas. (b) $\frac12\sqrt{186} = 6.8$ años y $\frac12\sqrt{2574} = 25.4$ años.
(c) $93/6.8 = 14$ y $1287/25.4 = 51$. El cociente entre $\frac12 n$ y $\frac12\sqrt{n}$ es
$\sqrt{n}$: crece como la raíz cuadrada del número de capas dudosas, así que la elección entre
las dos reglas importa cada vez más con la profundidad.
\end{solution}

\begin{solution}{exo:b2:climate-proxies:10}
(a) Sobre la \omterm{prop:b2:isotope-geochemistry:mwl}{recta de agua meteórica}, $\delta$D cambia ocho veces más que $\delta^{18}$O, así que la
pendiente es $8 \times 0.8 = 6.4$ por mil de $\delta$D por grado. La caída es $441.7 - 396.6 = 45.1$
por mil, y por tanto un enfriamiento de $45.1/6.4 = \qty{7.0}{\celsius}$. (b) El océano glaciar era unos
$8 \times 1.0 = 8$ por mil más pesado en $\delta$D; la nieve formada a partir del océano actual habría sido
más ligera en esa misma cantidad, $-441.7 - 8 = -449.7$ por mil. La caída pasa a ser de $53.1$ por mil y el enfriamiento
de $53.1/6.4 = \qty{8.3}{\celsius}$. (c) $9.39 - 0.09 = \qty{9.30}{\celsius}$, así que $45.1/9.30 = 4.85$
por mil por grado. La serie publicada también está corregida de la altitud cambiante del
casquete, y su relación con la temperatura procede de modelos que simulan los isótopos del agua, no solo de
la pendiente espacial.
\end{solution}

\begin{solution}{exo:b2:climate-proxies:11}
(a) Los logaritmos de las razones son $-0.528$, $-0.223$, $0.191$, $0.501$, $0.904$, $1.227$; su
media es $0.345$ y la temperatura media 14. Con $S_{TT} = \sum (T - 14)^2 = 280$ y $S_{T\ell} =
\sum (T - 14)(\ell - 0.345) = 24.93$, $A = 24.93/280 = \qty{0.0890}{\per\celsius}$ y $\ln B =
0.345 - 0.0890 \times 14 = -0.901$, así que $B = \qty{0.406}{mmol/mol}$. (b) Los residuos son
$+0.017$, $-0.034$, $+0.023$, $-0.023$, $+0.025$ y $-0.009$, unos pocos por ciento de la razón, sin
tendencia; el ajuste está cerca de $0.4\,\mathrm{e}^{0.09T}$. (c) $T = (\ln 1.30 + 0.901)/0.0890 = (0.262 +
0.901)/0.0890 = \qty{13.1}{\celsius}$. (d) $T = (\ln 0.45 + 0.901)/0.0890 = (-0.799 +
0.901)/0.0890 = \qty{1.2}{\celsius}$. Este valor queda por debajo de la muestra más fría de la calibración,
\qty{4}{\celsius}: la ley se extrapola, y se sabe que a bajas temperaturas la razón Mg/Ca de los caparazones
se aparta de la ley exponencial por influencia del ion carbonato.
\end{solution}

\begin{solution}{exo:b2:climate-proxies:12}
(a) Por unidad de área oceánica, el agua del océano actual, $H$ con 0, se reparte entre los depósitos
$h_i$ con $\delta_i$ y el océano restante $H - h$, con $h = \sum h_i$: $0 = \sum h_i\delta_i +
(H - h)\,\Delta\delta$. Tras la primera etapa, $\Delta\delta = 60 \times 30/(3682.2 - 60) =
1800/3622.2 = 0.50$ por mil; tras la segunda, $(1800 + 70 \times 45)/(3682.2 - 130) =
4950/3552.2 = 1.39$ por mil. (b) $-4950/130 = -38.1$ por mil. (c) La regla es $1.39/130 = 0.0107$
por mil por metro, y da $0.50/0.0107 = \qty{47}{m}$ para la primera etapa en lugar de
\qty{60}{m}: el primer hielo era más pesado que la media, así que cada uno de sus metros desplazó el océano
menos. Una regla lineal calibrada en un estado lee mal los demás cuando cambia la composición media del
hielo.
\end{solution}

\begin{solution}{pb:b2:climate-proxies:1}
\textbf{1.} $361.84 \times 10^6\ \unit{km^2} \times \qty{0.130}{km} = \qty{4.70e7}{km^3}$, unos 47
millones de kilómetros cúbicos.
\textbf{2.} $130 \times 40/(3682.2 - 130) = 5200/3552.2 = 1.46$ por mil.
\textbf{3.} $5000/3557.2 = 1.41$ y $5360/3548.2 = 1.51$ por mil.
\textbf{4.} $4550/3552.2 = 1.28$ y $6500/3552.2 = 1.83$ por mil. La composición del hielo
mueve la respuesta en \num{0.55} por mil, el nivel del mar en \num{0.10}: el hielo importa más.
\textbf{5.} $\delta_{\mathrm{ice}} = -1.0 \times 3552.2/130 = -27.3$ por mil.
\textbf{6.} $40/3682.2 = 0.0109$ por mil por metro, que se redondea a los \num{0.011} de los niveles
del mar datados.
\textbf{7.} Con $\delta^{18}\mathrm{O}_{\mathrm{sw}}$ constante, $\mathrm{d}T = -4.76\,\mathrm{d}\delta^{18}
\mathrm{O}_{\mathrm{b}}$, así que $\mathrm{d}T = -1$ da $\mathrm{d}\delta^{18}\mathrm{O}_{\mathrm{b}} =
1/4.76 = 0.21$ por mil.
\textbf{8.} $\delta^{18}\mathrm{O}_{\mathrm{b}} = (16.1 - T)/4.76 + \delta^{18}\mathrm{O}_{\mathrm{sw}}
- 0.27 = 14.1/4.76 - 0.27 = 2.96 - 0.27 = 2.69$ por mil.
\textbf{9.} $17.1/4.76 + 1.0 - 0.27 = 3.59 + 0.73 = 4.32$ por mil. El cambio es $4.32 - 2.69 =
1.63$ por mil: \num{1.00} debido al agua y $3.0 \times 0.21 = 0.63$ debido al enfriamiento de
\qty{3}{\celsius}.
\textbf{10.} $T = \ln(2.0/0.4)/0.09 = \ln 5/0.09 = \qty{17.9}{\celsius}$.
\textbf{11.} $0.18/0.1 = \qty{1.8}{\celsius}$ de enfriamiento.
\textbf{12.} La razón Mg/Ca es un cociente entre dos elementos de la calcita, fijado por la temperatura de
calcificación; no hace intervenir los isótopos del oxígeno, y por eso ignora la composición del
agua que cambia el hielo.
\textbf{13.} $4.92 - 3.23 = 1.69$ por mil.
\textbf{14.} Al resolver $\Delta = h \times 40/(H - h)$ respecto de $h$: $h = H\Delta/(\Delta + 40) = 3682.2
\times 1.69/41.69 = \qty{149}{m}$, más que los 125 a \qty{134}{m} de los indicios del nivel del mar:
parte del cambio tiene que ser enfriamiento.
\textbf{15.} $1.69 - 1.46 = 0.23$ por mil, y $4.76 \times 0.23 = \qty{1.1}{\celsius}$.
\textbf{16.} $1.69 - 1.00 = 0.69$ por mil, y $4.76 \times 0.69 = \qty{3.3}{\celsius}$ (3.28).
\textbf{17.} $0.23/0.25 = \qty{0.9}{\celsius}$ y $0.69/0.25 = \qty{2.8}{\celsius}$.
\textbf{18.} El Agua Profunda del Atlántico Norte se enfrió de 2--\qty{4}{\celsius} a \qty{-1.1}{\celsius},
es decir, entre 3.1 y \qty{5.1}{\celsius}; el Agua Antártica de Fondo, de 0--\qty{1}{\celsius} a
\qty{-1.9}{\celsius}, entre 1.9 y \qty{2.9}{\celsius}; la razón Mg/Ca del Atlántico noroccidental da
\qty{1.8}{\celsius}. Todos quedan entre unos 2 y \qty{5}{\celsius}, en torno a los 2.8 a
\qty{3.3}{\celsius} de la corrección del agua intersticial, y por encima del \qty{1}{\celsius} del balance
de masas con hielo a $-40$ por mil.
\textbf{19.} $-396.6 - (-441.7) = 45.1$ por mil.
\textbf{20.} $8 \times 0.8 = 6.4$ por mil por grado, así que $45.1/6.4 = \qty{7.0}{\celsius}$.
\textbf{21.} $9.39 - 0.09 = \qty{9.30}{\celsius}$, y $45.1/9.30 = 4.85$ por mil por grado.
\textbf{22.} Dome C, \qty{9.3}{\celsius}, más que la superficie global, de 5 a \qty{7}{\celsius}, más
que el océano profundo, unos \qty{3}{\celsius}.
\textbf{23.} El agua profunda ya está fría hoy, de 0 a \qty{4}{\celsius}, y el agua de mar no puede
enfriarse mucho por debajo de su punto de congelación, que el agua de fondo glaciar casi alcanzó a
\qty{-1.9}{\celsius}: el océano profundo solo tenía unos pocos grados que perder, mientras que el aire sobre la Antártida
no tiene ese suelo.
\textbf{24.} Con la corrección del agua intersticial, el océano profundo estaba unos \textbf{\qty{3.3}{\celsius}
más frío} en el Último Máximo Glacial que hoy, cerca de los 2 a \qty{5}{\celsius} de las
estimaciones independientes del agua profunda.
\end{solution}
